% 前言
\clearpage
\pagenumbering{arabic}
\setcounter{page}{1}

\chapter*{前\ \ 言}
\addcontentsline{toc}{chapter}{前言}
\markboth{前言}{前言}

生产实习是本科阶段将课堂所学理论知识与实际工程问题相结合的重要实践环节。近年来，深度生成模型在图像、语音等领域取得了显著进展，其中以去噪扩散概率模型（Denoising Diffusion Probabilistic Model）与基于分数的生成模型（Score-based Generative Model, SGM）为代表的 score-based 方法，凭借其训练稳定、生成质量高等优点，已成为生成建模领域的主流范式之一。与图像、语音等数据不同，工业场景中的多元时序数据（如设备传感器读数、环境监测数据、卫星遥测数据等）普遍存在标注稀缺、异常样本稀少、跨设备/跨工况分布差异大等问题，导致基于时序数据训练的异常检测、健康预测、故障诊断等模型往往面临数据不足的困境。将 score-based 生成模型引入时序数据合成，为缓解上述数据稀缺问题、开展数据增强与仿真验证提供了一条可行路径。

本课题选取了 AAAI 2023 论文《Regular Time-series Generation using SGM》（简称 TSGM）作为复现对象。该论文是较早将 score-based 生成模型系统性地应用于“规则时序生成”任务的工作，与已有的时序预测（如 TimeGrad、ScoreGrad）、时序插补（如 CSDI）等 score-based 时序任务相区别，其核心贡献在于提出了隐空间中的条件去噪分数匹配目标与递归式的 Predictor-Corrector 采样流程。本课题在原论文实验设置的基础上，将方法迁移到三个工业及环境监测相关的时序数据集（C-MAPSS 涡扇发动机退化仿真数据、PM25 空气质量数据、ESA 卫星遥测异常数据）上进行了完整复现，覆盖论文精读、数据适配、模型实现、两阶段训练、递归采样与标准评估协议等全流程环节，并对复现过程中遇到的工程问题与方法本身的局限性进行了具体分析。

本报告的撰写目的，一方面是对本次生产实习课程设计的工作内容、实现方案与实验结果进行系统总结；另一方面也希望通过对论文方法的深入复现与结果分析，加深对 score-based 生成模型在时序数据领域应用的理解。由于本人水平有限，报告中难免存在不足之处，恳请各位老师批评指正。
